\documentclass[11pt]{article}

\usepackage[margin=1in]{geometry}
\usepackage{amsmath, amssymb}
\usepackage{booktabs}
\usepackage{graphicx}
\usepackage{hyperref}
\usepackage{algorithm}
\usepackage{algpseudocode}
\usepackage{enumitem}
\usepackage{xcolor}

\hypersetup{colorlinks=true, linkcolor=blue, urlcolor=blue, citecolor=blue}

\title{A Markov-Chain Multi-Market Trading Strategy: \\
Methodology, Implementation, and Backtest Design}
\author{}
\date{\today}

\begin{document}
\maketitle

\begin{abstract}
This note documents a discrete-state Markov-chain trading strategy that is
applied across five global equity index universes (S\&P 500, Hang Seng,
Shanghai Composite, FTSE 100, Nikkei 225) and a fixed global macro basket
(crypto, commodities, and leveraged/sector ETFs). Daily returns are
z-scored and mapped onto a five-state alphabet, a train-only transition
matrix is estimated per ticker, and tickers are retained only if their
one-step-ahead directional signal clears a minimum edge and hit-rate
threshold on the training data. Retained tickers generate Kelly-capped,
fixed-holding-period trade signals on held-out test data, which are then
run through a daily event-driven backtest with transaction costs, intraday
stop losses, and a volatility/trend-based index-rotation fallback for days
without an active signal. We describe the full pipeline, the risk
management rules, the performance attribution methodology (Sharpe ratio
and CAPM alpha/beta against SPY), and discuss the backtest's principal
limitations, most notably survivorship bias in the ticker universes.
\end{abstract}

\tableofcontents

% ============================================================
\section{Motivation}
% ============================================================

Momentum and mean-reversion effects in daily equity returns are well
documented empirically, but their sign and strength are regime- and
asset-dependent. Rather than assume a single global rule (e.g. ``buy after
two down days''), this project estimates a \emph{per-ticker}, data-driven
transition structure: each ticker's own history is discretized into a small
state alphabet describing the size of its daily move, and a transition
matrix over that alphabet is fit on a training sample. If the fitted matrix
implies a strong enough directional edge out-of-sample, the ticker is
included in the tradable universe; if not, it is discarded. This turns
"does momentum/mean-reversion exist here" into an explicit, per-ticker
statistical test with a pre-registered threshold, rather than a strategy
applied uniformly to every instrument.

% ============================================================
\section{State Space and Transition Matrix Estimation}
% ============================================================

\subsection{Return discretization}

For a given ticker, let $r_t$ denote the daily percentage return computed
from adjusted close prices. The full return history is split
chronologically into a training segment (the first 70\% of observations)
and a held-out test segment (the remaining 30\%):
\[
  r_1, \dots, r_T \;=\; \underbrace{r_1, \dots, r_{\lfloor 0.7T \rfloor}}_{\text{train}}
  \;\cup\; \underbrace{r_{\lfloor 0.7T \rfloor + 1}, \dots, r_T}_{\text{test}}.
\]

The training-set mean $\mu$ and (population) standard deviation $\sigma$
are computed \emph{only} from the training segment, and used to z-score
\emph{both} segments:
\[
  z_t = \frac{r_t - \mu}{\sigma}, \qquad t = 1, \dots, T.
\]

Using train-only statistics to standardize the test segment avoids
look-ahead bias: nothing about the test-period distribution of returns
influences the discretization boundaries.

Each $z_t$ is then mapped to one of five states using two boundary
parameters, $\delta = 0.55$ (the ``moderate move'' boundary) and
$\varepsilon = 1.15$ (the ``extreme move'' boundary):

\[
  s_t =
  \begin{cases}
    \text{Very Down} & z_t \le -\varepsilon \\
    \text{Down}      & -\varepsilon < z_t \le -\delta \\
    \text{Neutral}   & -\delta < z_t < \delta \\
    \text{Up}        & \delta \le z_t < \varepsilon \\
    \text{Very Up}   & z_t \ge \varepsilon
  \end{cases}
\]

\subsection{Transition matrix estimation with Laplace smoothing}

Let $\mathcal{S} = \{\text{Very Down}, \text{Down}, \text{Neutral},
\text{Up}, \text{Very Up}\}$. From the training-set state sequence
$s_1, \dots, s_n$, raw transition counts are accumulated:
\[
  N_{ij} = \sum_{t=1}^{n-1} \mathbf{1}\{s_t = i,\, s_{t+1} = j\}, \qquad i, j \in \mathcal{S}.
\]

To avoid zero-probability transitions for state pairs that are rare or
absent in the training window, Laplace (additive) smoothing with constant
$\alpha = 0.7$ is applied before row-normalizing:
\[
  \hat{P}(j \mid i) = \frac{N_{ij} + \alpha}{\sum_{k \in \mathcal{S}} (N_{ik} + \alpha)}.
\]

This yields a $5 \times 5$ row-stochastic transition matrix $\hat{P}$,
estimated entirely from the training segment.

% ============================================================
\section{Edge Filtering}
% ============================================================

Not every ticker's transition matrix implies a usable trading edge. For
each state $i$, define the \emph{net directional probability}
\[
  \Delta(i) = \big[\hat P(\text{Up}\mid i) + \hat P(\text{Very Up}\mid i)\big]
            - \big[\hat P(\text{Down}\mid i) + \hat P(\text{Very Down}\mid i)\big].
\]

A ticker is tested for edge by walking the training-set state sequence: at
each step $t$ where $|\Delta(s_t)| \ge \tau$ (edge threshold, $\tau = 0.2$),
a directional prediction is made,
$\text{sign}(\Delta(s_t)) \in \{+1, -1\}$, and scored against the
\emph{realized} next-day return $r_{t+1}$ as a hit if
$\text{sign}(\Delta(s_t)) \cdot r_{t+1} > 0$.

Let $H$ be the number of such qualifying predictions and $h$ the number of
hits. The ticker clears the edge filter only if:
\begin{itemize}[nosep]
  \item $H \ge 50$ (minimum sample-size guard against spurious small-sample edges), and
  \item $h / H \ge 0.55$ (minimum hit rate).
\end{itemize}

Tickers that fail either condition are dropped from the tradable universe
entirely — no signals are generated from them on the test set. This is an
in-sample (training-set) filter; whether the edge \emph{persists}
out-of-sample is exactly what the test-set backtest evaluates.

% ============================================================
\section{Signal Generation and Position Sizing}
% ============================================================

For each ticker that clears the edge filter, its transition matrix
$\hat P$ is applied to the \textbf{held-out test-set} state sequence to
generate trade signals. At test-set time $t$ with state $s_t$:
\[
  \Delta_{\text{long}} = \max(\Delta(s_t), 0), \qquad
  \Delta_{\text{short}} = \max(-\Delta(s_t), 0).
\]

A signal fires only if $\max(\Delta_{\text{long}}, \Delta_{\text{short}})
\ge \tau$. The trade direction is Long if $\Delta_{\text{long}} \ge
\Delta_{\text{short}}$, else Short. The position's raw Kelly-style weight
is
\[
  \kappa = \min\big(\max(\Delta_{\text{long}}, \Delta_{\text{short}}),\; \kappa_{\max}\big),
  \qquad \kappa_{\max} = 0.20,
\]
i.e. the magnitude of the directional edge itself is used as a
(capped) sizing signal — larger edges get proportionally larger size, up to
a 20\% cap. Each signal has a fixed holding period of $H_d = 5$ trading
days: the position is entered at the test-set date's close and scheduled to
exit five trading days later.

\subsection{Daily capital allocation}

On any trading day $d$ with one or more fresh signals, a fixed fraction of
the day's starting equity is put at risk:
\[
  C_d = \phi \cdot E_d, \qquad \phi = 0.55,
\]
where $E_d$ is equity (cash plus mark-to-market of open positions) at the
start of day $d$. At most $M = 3$ signals are taken per day, chosen by
descending $\kappa$, and each is allocated a share of $C_d$ proportional to
its own $\kappa$ relative to the sum of $\kappa$ across the day's taken
signals.

% ============================================================
\section{Risk Management}
% ============================================================

\subsection{Intraday stop losses}

Every open position (other than on its entry day) is checked daily against
an intraday stop, referenced to the position's \emph{entry-day opening
price} $O_0$ (not the entry fill price):
\[
  \text{Long stop: } L_d \le 0.97 \cdot O_0,
  \qquad
  \text{Short stop: } H_d \ge 1.03 \cdot O_0,
\]
where $L_d, H_d$ are the day's low and high. If triggered, the position is
closed at a fixed booked loss of $\text{STOP\_LOSS\_PCT} = -3\%$ of its
notional size (rather than the exact intraday execution price), net of a
half round-trip transaction cost.

\subsection{Transaction costs}

A flat round-trip cost of 20 basis points ($\text{ROUND\_TRIP\_COST} =
0.002$) is charged on both entry and exit, applied as half the round-trip
cost against notional size at each leg.

\subsection{Quiet-day index rotation fallback}

On days with no qualifying Markov signal, the strategy does not sit
entirely idle. Instead, it ranks a fixed basket of eleven broad index ETFs
(SPY, QQQ, IWM, EFA, EEM, EWJ, FXI, EWU, EWG, EWQ, EWZ) by a
volatility-adjusted trend score:
\[
  \text{score}(i) = \frac{\hat\beta_i \cdot 252}{\hat\sigma_i},
\]
where $\hat\beta_i$ is the slope of a linear fit to $\log$-price over the
trailing 60 trading days, and $\hat\sigma_i$ is the trailing 20-day
realized volatility, annualized. The strategy goes long the
highest-scoring ETF and short the lowest-scoring ETF, sized inversely to
their respective volatilities and capped at 25\% of equity in aggregate
($\text{FALLBACK\_BUDGET\_FRAC}$), held for 5 trading days.

% ============================================================
\section{Backtest Engine}
% ============================================================

The simulation is a daily event loop over the full business-day range
implied by the generated signals. Algorithm~\ref{alg:loop} summarizes the
per-day sequence of operations.

\begin{algorithm}[h]
\caption{Daily backtest loop}
\label{alg:loop}
\begin{algorithmic}[1]
\For{each business day $d$}
  \State Mark all open positions to market; compute $E_d$ = cash + MTM
  \If{signals exist for $d$}
    \State Open Markov-signal trades, sized per Section 4.1
  \Else
    \State Open index-rotation fallback trades (Section 5.3)
  \EndIf
  \State Apply intraday stop losses to all positions opened on prior days
  \State Close any position whose scheduled exit day has arrived
  \State Record end-of-day equity and daily P\&L
\EndFor
\end{algorithmic}
\end{algorithm}

Price data is cached in two forms prior to simulation: an
adjusted-close series (for mark-to-market and exit pricing) and a
\emph{non-adjusted} OHLC series (needed so that intraday stop checks
reference true traded highs/lows rather than dividend/split-adjusted
values).

% ============================================================
\section{Performance Attribution}
% ============================================================

Let $\{E_d\}$ be the realized daily equity curve and $\{\rho_d\}$ the
corresponding simple daily returns. With an assumed annual risk-free rate
of $5\%$ ($\rho_f = (1.05)^{1/252} - 1$ daily), the excess return series is
$\rho_d^{\text{ex}} = \rho_d - \rho_f$.

\paragraph{Sharpe ratio.}
\[
  \text{Sharpe}_{\text{daily}} = \frac{\overline{\rho^{\text{ex}}}}{\text{sd}(\rho^{\text{ex}})},
  \qquad
  \text{Sharpe}_{\text{annual}} = \sqrt{252}\cdot \text{Sharpe}_{\text{daily}}.
\]

\paragraph{Annualized return.} Both a compounding-based estimate and a
calendar-time CAGR are reported:
\[
  r_{\text{ann}} = \Big(\prod_d (1+\rho_d)\Big)^{252/n} - 1,
  \qquad
  \text{CAGR} = \Big(\frac{E_{\text{final}}}{E_{\text{initial}}}\Big)^{365.25/\text{days}} - 1.
\]

\paragraph{CAPM alpha and beta.} The strategy's excess returns are
regressed on SPY's excess returns over the same aligned dates:
\[
  \rho_d^{\text{ex}} = \alpha + \beta \cdot \rho_{d,\text{SPY}}^{\text{ex}} + \epsilon_d,
\]
estimated by OLS. Both a linearly annualized alpha ($252\alpha$) and a
compounded annualized alpha ($(1+\alpha)^{252}-1$) are reported, along with
$\beta$, $R^2$, and the $t$-statistic/$p$-value on $\alpha$ for a rough
read on statistical significance of any outperformance.

% ============================================================
\section{Discussion and Limitations}
% ============================================================

\begin{itemize}
  \item \textbf{Survivorship bias.} Index constituents (S\&P 500, FTSE 100,
  Hang Seng, Nikkei 225) are pulled from \emph{current} membership lists
  and then backtested over 20 years of price history. Firms removed from
  an index over that period (delisted, acquired, bankrupt) are excluded
  from the universe, which mechanically biases the backtest's historical
  performance upward relative to a strategy that could only have known
  index membership in real time.
  \item \textbf{In-sample edge filter, out-of-sample evaluation.} The edge
  filter (Section 3) is applied to the training segment; the actual P\&L
  is generated only from test-segment signals. This mitigates but does not
  eliminate selection effects from testing many tickers simultaneously
  (multiple-comparisons risk): with a large universe, some tickers will
  clear a 55\% training-set hit-rate threshold by chance alone.
  \item \textbf{Approximated stop-loss execution.} A triggered stop books a
  fixed $-3\%$ loss on notional rather than the exact price at which a
  gap-through would have executed, which understates losses in the tails
  (large overnight gaps).
  \item \textbf{No cross-market FX adjustment.} Hang Seng, Shanghai, FTSE,
  and Nikkei tickers are priced and P\&L'd in local currency; no
  currency-hedging or FX-translation is modeled, so realized USD returns
  for a US-based investor would differ from what is reported.
  \item \textbf{Data coverage and runtime.} Data availability and quality
  via \texttt{yfinance} varies substantially across exchanges; the
  volatility-based prefilter (retaining only the top 20\% most volatile
  tickers) is a practical runtime mitigation, not a methodological
  requirement of the strategy.
\end{itemize}

This document accompanies the code in the project repository
(\texttt{markov\_strategy/}) and is intended as a methodology reference,
not as investment advice. Past backtest performance is not indicative of
future results.

\end{document}
